% ============================================================
% Capitolo 2 — Setup del primo progetto in 15 minuti
% ============================================================

\chapter{Setup del primo progetto in 15 minuti}
\label{ch:setup-primo-progetto}

Lorenzo ha clonato il repository di TaskFlow API, aperto il terminale e ha davanti a sé una directory con il codice del progetto e nient'altro. Nessun \texttt{\_CONTEXT.md}, nessun framework AI. Questo è il punto di partenza.

In questo capitolo seguiamo Lorenzo --- e tu con lui --- mentre aggiunge AI-DLC al progetto, compila il contesto per la prima volta e fa il primo turno con un agente AI. Quindici minuti, non di più.

\section{Passo 1 --- Copiare il framework nel repository}
\label{sec:02-passo1-copia}

AI-DLC non ha un installer. Si copia. La struttura che ti serve è questa:

\begin{lstlisting}[language=,caption={Struttura del framework nel repository}]
progetto/
+-- AGENTS.md          ← OpenAI Codex CLI
+-- CLAUDE.md          ← Claude Code
+-- GEMINI.md          ← Google Gemini
+-- OPENCLAW.md        ← OpenClaw
+-- VISUALSTUDIO.md    ← Visual Studio 2026
+-- .github/
|   +-- copilot-instructions.md  ← GitHub Copilot
+-- .ai-dlc/
    +-- halt-triggers.yaml
    +-- manifest.json
    +-- modules/       ← regole del framework
    +-- schemas/       ← schemi JSON
    +-- tools/         ← script di utilità
\end{lstlisting}

\begin{notebox}
L'URL del repository è indicativo --- verifica quello aggiornato nella documentazione ufficiale del framework o nel file \texttt{README.md} della distribuzione che hai ricevuto.
\end{notebox}

Clona il framework in una directory temporanea e copia i file:

\begin{lstlisting}[language=bash,caption={Installazione su Bash (Linux/macOS/WSL)}]
git clone https://github.com/rollaradio/adlc-framework /tmp/adlc
cp /tmp/adlc/AGENTS.md .
cp /tmp/adlc/CLAUDE.md .
cp -r /tmp/adlc/.github .
cp -r /tmp/adlc/.ai-dlc .
\end{lstlisting}

\begin{lstlisting}[language=,caption={Installazione su PowerShell (Windows)}]
git clone https://github.com/rollaradio/adlc-framework $env:TEMP\adlc
Copy-Item "$env:TEMP\adlc\AGENTS.md" .
Copy-Item "$env:TEMP\adlc\CLAUDE.md" .
Copy-Item "$env:TEMP\adlc\.github" . -Recurse
Copy-Item "$env:TEMP\adlc\.ai-dlc" . -Recurse
\end{lstlisting}

\begin{notebox}
I file \texttt{AGENTS.md}, \texttt{CLAUDE.md}, ecc. sono read-only per l'uso normale --- non modificarli. Le tue personalizzazioni vanno in \texttt{.ai-dlc/project/} (lo vediamo al §\ref{sec:02-project-dir}).
\end{notebox}

\section{Passo 2 --- Scaffold dello stato di progetto}
\label{sec:02-passo2-scaffold}

Esegui lo script di inizializzazione per creare i file di stato:

\begin{lstlisting}[language=bash,caption={Inizializzazione su Bash}]
bash .ai-dlc/tools/init.sh
\end{lstlisting}

\begin{lstlisting}[language=,caption={Inizializzazione su PowerShell}]
.\.ai-dlc\tools\init.ps1
\end{lstlisting}

Lo script crea quattro cose:

\begin{table}[htbp]
  \centering
  \begin{tabular}{ll}
    \toprule
    \textbf{File/Cartella} & \textbf{A cosa serve} \\
    \midrule
    \texttt{\_CONTEXT.md} & Stato corrente (fase, task, vincoli) \\
    \texttt{PROGRESS.md} & Diario di bordo delle sessioni \\
    \texttt{.ai-dlc/project/instructions.md} & Regole specifiche per il progetto \\
    \texttt{.ai-dlc/project/skills/} & Skill personalizzate \\
    \bottomrule
  \end{tabular}
  \caption{File creati dallo script di inizializzazione}
  \label{tab:02-init-files}
\end{table}

Per progetti piccoli o spike usa il context minimo:

\begin{lstlisting}[language=bash,caption={Inizializzazione con context minimale}]
bash .ai-dlc/tools/init.sh --context minimal
\end{lstlisting}

\section{Passo 3 --- Compilare \texttt{\_CONTEXT.md}}
\label{sec:02-passo3-context}

Apri \texttt{\_CONTEXT.md} e compila i campi con i valori del tuo progetto. Ecco quello di Lorenzo per TaskFlow API:

\begin{lstlisting}[language=,caption={\_CONTEXT.md di TaskFlow API --- setup iniziale}]
| Param                   | Value                                    |
|-------------------------|------------------------------------------|
| Project                 | TaskFlow API                             |
| Phase                   | 0-Discovery                              |
| Mode                    | STANDARD                                 |
| Active Task             | T-001 Setup e analisi requisiti          |
| Active Task Token Est.  | 8000/2000/10000                          |
| Active Task Model Level | 3 Sonnet Low                             |

## Stack
| Layer     | Technology              |
|-----------|-------------------------|
| Runtime   | Node.js 22              |
| Framework | Fastify 5               |
| Database  | PostgreSQL 16 + Prisma  |
| Auth      | OIDC + JWT              |

## Security Constraints
| SEC-01 | Input Validation | Zod schema su ogni endpoint     |
| SEC-02 | Authentication   | OIDC, JWT con refresh rotation  |
| SEC-03 | Logging          | Nessun token o PII nei log      |

## Performance Constraints
| PERF-01 | Latency     | P95 < 200ms per endpoint |
| PERF-02 | Payload size| Max 1MB per risposta     |
\end{lstlisting}

\begin{tipbox}
Compila almeno \texttt{Project}, \texttt{Phase}, \texttt{Mode} e \texttt{Stack}. I vincoli SEC e PERF sono fondamentali --- anche solo due o tre, quelli più rilevanti. Il resto puoi aggiungerlo progressivamente.
\end{tipbox}

\section{Passo 4 --- Verificare con il validator}
\label{sec:02-passo4-validate}

\begin{lstlisting}[language=bash,caption={Validazione su Bash}]
bash .ai-dlc/tools/validate.sh
\end{lstlisting}

\begin{lstlisting}[language=,caption={Validazione su PowerShell}]
.\.ai-dlc\tools\validate.ps1
\end{lstlisting}

Output atteso:

\begin{lstlisting}[language=,caption={Output del validator}]
AI-DLC validation — v4.0.0

[OK] Framework files present
[OK] AGENTS.md, CLAUDE.md, GEMINI.md, OPENCLAW.md, VISUALSTUDIO.md found
[OK] halt-triggers.yaml valid
[OK] manifest.json valid
[!]  _CONTEXT.md: Active Task Token Est. is placeholder (0/0/0)

AI-DLC validation passed (2 warnings)
\end{lstlisting}

I warning sui placeholder sono normali per la prima sessione. Con \texttt{--strict} diventano failure (utile in CI).

\section{Passo 5 --- Il primo turno con l'agente}
\label{sec:02-passo5-bootstrap}

Apri l'agente AI e chiedi la conferma del contesto:

\begin{quote}
``Carica \texttt{\_CONTEXT.md} e confermami il contesto: fase, task attivo, stack, vincoli SEC e PERF attivi.''
\end{quote}

L'agente risponde con qualcosa del tipo:

\begin{lstlisting}[language=,caption={Risposta di bootstrap dell'agente}]
Context: TaskFlow API | Phase: 0-Discovery | Task: T-001 |
Stack: Node.js/Fastify/PostgreSQL |
Constraints: SEC-01,SEC-02,SEC-03,PERF-01. Proceed?
\end{lstlisting}

Quando vedi questa conferma, sai che l'agente ha letto il contesto. Da questo momento puoi iniziare il dialogo reale senza ripetere nulla.

\section{Struttura raccomandata per \texttt{.ai-dlc/project/}}
\label{sec:02-project-dir}

\begin{lstlisting}[language=,caption={Struttura della cartella di personalizzazione}]
.ai-dlc/project/
+-- instructions.md   ← regole specifiche del progetto
+-- domain.md         ← business logic e regole di dominio
+-- conventions.md    ← convenzioni di codice
+-- skills/           ← skill personalizzate
\end{lstlisting}

Per TaskFlow API, Lorenzo aggiunge a \texttt{instructions.md}:

\begin{lstlisting}[language=,caption={.ai-dlc/project/instructions.md di TaskFlow API}]
# TaskFlow API — Project Instructions

## Conventions
- All API errors must use RFC 7807 Problem Details format
- Database queries use Prisma only — no raw SQL except in migrations
- All endpoints must have a corresponding integration test

## Git workflow
- Branch naming: feat/T-NNN-description, fix/T-NNN-description
- PR review required before merge to main
- Squash merge only
\end{lstlisting}

\section*{\LabelSummary}

\begin{enumerate}
  \item Copia i file del framework nella radice del progetto
  \item Esegui \texttt{init.ps1} / \texttt{init.sh} per lo scaffold iniziale
  \item Compila \texttt{\_CONTEXT.md} con fase, stack e vincoli
  \item Verifica con \texttt{validate.ps1} / \texttt{validate.sh}
  \item Fai il primo bootstrap con l'agente e ricevi conferma del contesto
\end{enumerate}

TaskFlow API è pronta. L'agente sa dove sei, cosa stai costruendo e cosa non può ignorare. Nel prossimo capitolo vediamo come si svolge una sessione completa, dall'inizio al commit finale.
